\chapter{Multi-Body Dynamics of Biological Chains}
\label{ch:multibody-bio}

\epigraph{%
  The human body is the most complex open kinematic chain in nature.
  Analyzing it with the tools of rigid-body dynamics is both
  necessary and insufficient.
}{}

\section{Open-Chain vs.\ Closed-Chain Biomechanics}

Most human movement involves open kinematic chains (the limbs) connected to a floating
base (the pelvis/trunk). However, many functional activities create temporary closed
chains: the stance phase of gait (foot on ground), pushing against a wall, or the golf
swing when both hands grip the club.

\subsection{Equations of Motion for Biological Chains}

The equations of motion for an $n$-DOF musculoskeletal system take the standard
Lagrangian form:
\begin{equation}
  M(\bm{q})\ddot{\bm{q}} + C(\bm{q}, \dot{\bm{q}})\dot{\bm{q}} + g(\bm{q})
  = R(\bm{q}) \bm{F}_{\text{muscle}} + J_c^T(\bm{q}) \bm{F}_{\text{contact}},
  \label{eq:ch5:eom-bio}
\end{equation}
where:
\begin{itemize}
  \item $M(\bm{q}) \in \Reals^{n \times n}$ is the mass matrix (from segment inertial properties)
  \item $C(\bm{q}, \dot{\bm{q}}) \in \Reals^{n \times n}$ is the Coriolis/centrifugal matrix
  \item $g(\bm{q}) \in \Reals^n$ is the gravity vector
  \item $R(\bm{q}) \in \Reals^{n \times p}$ is the muscle moment arm matrix ($p$ muscles, $n$ DOF)
  \item $\bm{F}_{\text{muscle}} \in \Reals^p$ are the muscle forces (from Hill models, Chapter~3)
  \item $J_c(\bm{q})$ is the contact Jacobian
  \item $\bm{F}_{\text{contact}}$ are external contact forces
\end{itemize}

The key difference from Volume~I's $\bm{\tau} = B\bm{u}$ is that joint torques are
no longer direct control inputs. Instead:
\begin{equation}
  \bm{\tau}_{\text{net}} = R(\bm{q}) \bm{F}_{\text{muscle}}
  = \sum_{i=1}^{p} \bm{r}_i(\bm{q}) \cdot F_{\text{muscle},i},
  \label{eq:ch5:net-torque}
\end{equation}
where $\bm{r}_i(\bm{q}) \in \Reals^n$ is the $i$-th column of $R$, giving the moment
arm vector of muscle $i$ across all joints.

\subsection{Bi-articular Muscles}

Unlike engineering actuators that span a single joint, many muscles cross two
or more joints. These \textbf{bi-articular} (or multi-articular) muscles create
mechanical coupling between joints that has no analog in standard robotics.

\begin{principle}{The Coordination Function of Bi-articular Muscles}{biarticular}
  Bi-articular muscles serve two distinct functions:
  \begin{enumerate}
    \item \textbf{Energy transfer}: they transport mechanical energy from one joint to
          another without metabolic cost. When the proximal joint extends (muscle
          shortening at proximal end) and the distal joint flexes (muscle lengthening
          at distal end), energy is transferred proximally to distally.
    \item \textbf{Coordination constraint}: they link joint motions, reducing the
          effective degrees of freedom and simplifying the control problem.
  \end{enumerate}
  This is a form of \emph{mechanical intelligence}---the morphology of the
  musculoskeletal system performs part of the control task. The gastrocnemius,
  for example, crosses both the knee and ankle, coupling knee extension to ankle
  plantarflexion during the push-off phase of gait.

  \noindent
  \textbf{Control-Theoretic Perspective:} In the language of Agrachev \& Sachkov's geometric
  control theory \cite{AgrachevSachkov2004}, the moment arm matrix $R(\bm{q})$ is a covector
  field describing how muscle forces map to joint torques. Bi-articular muscles induce dependencies
  in the columns of $R(\bm{q})$, effectively constraining the input distribution $G(\bm{x})$ of the
  system. This reduces the effective control dimension and is formalized in Chapter~3 of
  \cite{AgrachevSachkov2004} (reachability and controllability of affine systems).
\end{principle}

\begin{lstlisting}[language=Python, caption={Computing joint torques from muscle forces through moment arms.}]
import numpy as np

def compute_joint_torques(
    moment_arms: np.ndarray,
    muscle_forces: np.ndarray
) -> np.ndarray:
    """Compute net joint torques from muscle forces and moment arms.

    Parameters
    ----------
    moment_arms : np.ndarray, shape (n_joints, n_muscles)
        Moment arm matrix R(q). Each column gives the moment arm
        of one muscle across all joints. Sign convention: positive
        moment arm produces positive joint torque (typically flexion).
    muscle_forces : np.ndarray, shape (n_muscles,)
        Force produced by each muscle (always positive).

    Returns
    -------
    np.ndarray, shape (n_joints,)
        Net joint torques.
    """
    return moment_arms @ muscle_forces


# Example: 2-DOF elbow system with 4 muscles
# Joints: shoulder flexion, elbow flexion
# Muscles: anterior deltoid, biceps (biarticular),
#          triceps (biarticular), brachialis
#
# Convention: R has shape (n_joints, n_muscles) so that
# tau = R @ F maps the muscle-force vector to a joint-torque
# vector via a single matrix-vector product. Each row is a joint
# and each column is a muscle. This is the shape that
# compute_joint_torques() expects.
R = np.array([
    # deltoid, biceps, triceps, brachialis
    [0.05,  0.03,  -0.025,  0.0 ],  # row 0: shoulder moment arms (m)
    [0.0,   0.04,  -0.020,  0.03],  # row 1: elbow moment arms (m)
])  # Shape: (2 joints, 4 muscles) -- ready to use as-is

# Note: biceps has moment arms at BOTH joints (bi-articular)
muscle_forces = np.array([200.0, 150.0, 100.0, 120.0])  # Newtons

tau = compute_joint_torques(R, muscle_forces)
print(f"Shoulder torque: {tau[0]:.1f} N*m")
print(f"Elbow torque:    {tau[1]:.1f} N*m")
\end{lstlisting}

\subsection{Co-Contraction and Joint Stiffness}

When antagonist muscles co-contract (both agonist and antagonist are active
simultaneously), the net joint torque may be zero but the joint \textbf{stiffness}
increases. This is because each active muscle acts as a spring:
\begin{equation}
  K_{\text{joint}} = \sum_{i=1}^{p} r_i^2(\bm{q}) \cdot \frac{\partial F_{\text{muscle},i}}{\partial \ell_i^M},
  \label{eq:ch5:joint-stiffness}
\end{equation}
where $\partial F / \partial \ell^M$ is the short-range stiffness of the muscle, which
depends on activation level.

Co-contraction is metabolically expensive but biomechanically important: it allows the
nervous system to modulate joint \emph{impedance} independently of joint
\emph{torque}. This is a form of \textbf{variable impedance control} that has no direct
analog in rigid robots with torque-controlled joints.

Position-dependent stiffness modulation is a nonlinear phenomenon central to the geometric
analysis of affine control systems. Agrachev \& Sachkov \cite{AgrachevSachkov2004} develop
the theory of feedback linearization and stabilization via nonlinear drift terms (Chapter~6),
which encompasses exactly this scenario: using position-dependent actuation to achieve
desired closed-loop dynamics.

\section{The Mass Matrix of the Human Body}

Computing the mass matrix $M(\bm{q})$ for a full musculoskeletal model follows the
same recursive algorithms as Volume~0, Chapter~10 (the Composite Rigid Body Algorithm),
using the segment inertial properties from Chapter~2. The key computational steps are:

\begin{enumerate}
  \item Define the kinematic tree (segment connectivity and joint types)
  \item Compute forward kinematics using the product of exponentials
  \item Apply the CRBA to compute $M(\bm{q})$
\end{enumerate}

For a 20-DOF model (pelvis + 2 legs + trunk + 2 arms), $M(\bm{q})$ is a
$20 \times 20$ symmetric positive-definite matrix.

\section*{Exercises}

\begin{enumerate}
  \item \textbf{Bi-articular Energy Transfer.}
        For the gastrocnemius (crossing knee and ankle), show mathematically how
        knee extension can produce ankle torque through the bi-articular coupling.

  \item \textbf{Co-Contraction Analysis.}
        Compute the joint stiffness when biceps and triceps co-contract at 50\%
        activation with the elbow at $90^\circ$.

  \item \textbf{Mass Matrix Properties.}
        For a 2-segment arm model, compute $M(\bm{q})$ at three configurations and
        verify positive-definiteness. Plot the condition number as a function of elbow angle.

  \item \textbf{(Programming.)} Implement the CRBA for a 3-segment leg model using
        the de~Leva segment properties. Compare with Pinocchio output.

  \item \textbf{(Programming.)} Simulate a 2-DOF arm with 4 muscles (as in the code
        listing) performing a reaching movement. Use prescribed muscle activations
        and compute the resulting trajectory via forward dynamics.
\end{enumerate}
